\documentclass[11pt,a4paper]{article}
% Shared preamble for the robodog engineering documentation.
\usepackage[a4paper,margin=24mm,top=22mm,bottom=24mm]{geometry}
\usepackage[T1]{fontenc}
\usepackage[utf8]{inputenc}
\usepackage{lmodern}
\usepackage{microtype}
\usepackage{graphicx}
\usepackage{booktabs}
\usepackage{longtable}
\usepackage{array}
\usepackage{amsmath}
\usepackage{xcolor}
\usepackage{listings}
\usepackage{tcolorbox}
\usepackage{enumitem}
\usepackage{caption}
\usepackage{float}
\usepackage[hidelinks,bookmarks=true]{hyperref}

\definecolor{rdblue}{HTML}{1F4E79}
\definecolor{rdgrey}{HTML}{5A6675}
\definecolor{rdbg}{HTML}{F4F6F9}
\definecolor{rdline}{HTML}{C8D0DA}
\definecolor{rdred}{HTML}{A4262C}

\hypersetup{linkcolor=rdblue, citecolor=rdblue, urlcolor=rdblue,
            colorlinks=true, pdftitle={robodog -- software architecture}}

\captionsetup{font=small, labelfont={bf,color=rdblue}, width=0.92\textwidth}
\setlist{nosep, leftmargin=1.4em}
\renewcommand{\arraystretch}{1.15}

\lstset{
  basicstyle=\ttfamily\footnotesize,
  backgroundcolor=\color{rdbg},
  frame=single, rulecolor=\color{rdline},
  breaklines=true, showstringspaces=false,
  keywordstyle=\color{rdblue}\bfseries,
  commentstyle=\color{rdgrey}\itshape,
  numbers=none, xleftmargin=2pt, framexleftmargin=2pt,
  aboveskip=6pt, belowskip=6pt,
}

% Decision box: used for every choice that a reader might otherwise have to
% reverse-engineer from the code.
\newtcolorbox{decision}[1]{
  colback=rdbg, colframe=rdblue, boxrule=0.6pt, arc=2pt,
  left=6pt, right=6pt, top=5pt, bottom=5pt,
  fonttitle=\bfseries\small, title={Decision --- #1}}

\newtcolorbox{caveat}[1]{
  colback=white, colframe=rdred, boxrule=0.6pt, arc=2pt,
  left=6pt, right=6pt, top=5pt, bottom=5pt,
  fonttitle=\bfseries\small, title={#1}}

\newcommand{\unit}[1]{\,\mathrm{#1}}
\newcommand{\topicname}[1]{\texttt{\small #1}}
\newcommand{\code}[1]{\texttt{\small #1}}

% figure with a standard width and placement
\newcommand{\archfig}[3]{%
  \begin{figure}[H]\centering
  \includegraphics[width=#3\textwidth]{figures/#1.pdf}
  \caption{#2}\label{fig:#1}\end{figure}}

% AUTO-GENERATED by tools/make_doc_facts.py -- do not edit.
% Derived from robot_parameters.yaml (CAD 600e4a0e7d58f2e2)
% and robstride02.yaml. Regenerate after any model change.

\newcommand{\RdThighLength}{213}
\newcommand{\RdShankLength}{217.32}
\newcommand{\RdHaaX}{161}
\newcommand{\RdHaaY}{60}
\newcommand{\RdHfeDx}{60}
\newcommand{\RdHfeDr}{16}
\newcommand{\RdThighLat}{68.5}
\newcommand{\RdLateral}{84.5}
\newcommand{\RdReach}{430.32}
\newcommand{\RdFootRadius}{20}
\newcommand{\RdFootprintL}{442}
\newcommand{\RdFootprintW}{289}
\newcommand{\RdBodyL}{235}
\newcommand{\RdBodyW}{120}
\newcommand{\RdBodyH}{115}
\newcommand{\RdMassTotal}{10}
\newcommand{\RdMassStruct}{8.892}
\newcommand{\RdMassPayload}{1.04}
\newcommand{\RdMassCamera}{0.068}
\newcommand{\RdMassActuator}{380}
\newcommand{\RdMassTwelve}{4.56}
\newcommand{\RdMassBase}{3.934}
\newcommand{\RdMassHip}{624.5}
\newcommand{\RdMassThigh}{703.4}
\newcommand{\RdMassCalf}{171.7}
\newcommand{\RdComX}{-3.51}
\newcommand{\RdStandHeight}{320}
\newcommand{\RdStandHfe}{0.8098}
\newcommand{\RdStandKfe}{-1.5989}
\newcommand{\RdZeroHeight}{450.3}
\newcommand{\RdCrouchHeight}{240}
\newcommand{\RdRestHeight}{160}
\newcommand{\RdStandExtension}{69.7}
\newcommand{\RdStaticLoad}{24.52}
\newcommand{\RdHaaLo}{-0.8}
\newcommand{\RdHaaHi}{0.8}
\newcommand{\RdHfeLo}{-1.4}
\newcommand{\RdHfeHi}{2.4}
\newcommand{\RdKfeLo}{-2.6}
\newcommand{\RdKfeHi}{0}
\newcommand{\RdHaaDeg}{45.8}
\newcommand{\RdTorqueCont}{6}
\newcommand{\RdTorquePeak}{17}
\newcommand{\RdTorqueRated}{6}
\newcommand{\RdSpeedRated}{10.47}
\newcommand{\RdSpeedNoLoad}{42.94}
\newcommand{\RdSpeedOp}{20}
\newcommand{\RdGearRatio}{7.75}
\newcommand{\RdKt}{1.22}
\newcommand{\RdRphase}{0.29}
\newcommand{\RdVoltage}{48}
\newcommand{\RdVoltMin}{24}
\newcommand{\RdVoltMax}{60}
\newcommand{\RdEncoderBits}{14}
\newcommand{\RdEncoderRes}{\ensuremath{3.835\times 10^{-4}}}
\newcommand{\RdArmature}{\ensuremath{4.805\times 10^{-3}}}
\newcommand{\RdRotorInertia}{\ensuremath{8\times 10^{-5}}}
\newcommand{\RdTempWarn}{70}
\newcommand{\RdTempDerate}{85}
\newcommand{\RdTempFault}{100}
\newcommand{\RdRthermal}{2.84}
\newcommand{\RdCthermal}{190}
\newcommand{\RdTauThermal}{9}
\newcommand{\RdCanRate}{1}
\newcommand{\RdCanBuses}{2}
\newcommand{\RdControlRate}{400}
\newcommand{\RdBitsPerFrame}{150}
\newcommand{\RdBusLoad}{72}
\newcommand{\RdRatedArms}{4.95}
\newcommand{\RdCopperLoss}{21.3}
\newcommand{\RdCadSha}{600e4a0e7d58f2e2}


\title{\vspace{-12mm}\textbf{robodog}\\[2mm]
\large Software architecture for a 12-DOF quadruped\\
\normalsize 12 $\times$ ROBSTRIDE02 QDD actuators \textbullet\ \RdMassTotal\,kg \textbullet\ ROS 2 Jazzy}
\author{\small Generated from the Onshape CAD export \texttt{\RdCadSha}}
\date{\small Revision 0.1 --- simulation-complete, hardware-ready}

\begin{document}
\maketitle
\vspace{-6mm}

\begin{abstract}\noindent\small
This document specifies the first software architecture for a 12-DOF quadruped
built around twelve ROBSTRIDE02 quasi-direct-drive actuators. The robot model is
derived programmatically from the CAD assembly; the control stack runs at
\RdControlRate\,Hz behind a hardware abstraction that has three
implementations, two simulated and one for the real CAN bus. Everything is
implemented and running in simulation. Nothing above the hardware boundary
changes when the real actuators and camera are connected --- that is the
property the architecture is organised around, and Section~\ref{sec:migration}
states exactly what remains to be done.
\end{abstract}

\vspace{-2mm}
\tableofcontents
\vspace{4mm}

% =====================================================================
\section{Scope and status}

\paragraph{What exists.} Eight ROS 2 packages, 217 tests, a MuJoCo
model generated from the same parameters as the URDF, a five-room test house,
a modular web GUI, and the complete ROBSTRIDE02 CAN frame codec with 29 unit
tests. In simulation the robot stands, holds pose, and walks and trots at the
commanded velocity --- 0.30\,m/s tracked to within 3\,mm/s at 90\,\% of the
continuous torque rating.

\paragraph{What does not.} No hardware has been driven. The CAN backend is
written but unvalidated; per-joint direction and offset are unknown until the
robot is assembled and calibrated. Of the five gaits, pace and bound do not
track velocity --- neither is intended for travel, but neither is finished
either. Turning is implemented and untested beyond straight-line commands. The camera depth path is a placeholder in
one respect only (how the HP60C delivers depth over USB). Estimated quantities
--- rotor inertia, thermal constants, joint friction --- are marked as such
throughout and listed in Section~\ref{sec:assumptions}.

\paragraph{How to read this.} Every quantity quoted here is generated from
\code{robot\_parameters.yaml} and \code{robstride02.yaml} by
\code{tools/make\_doc\_facts.py}. If the CAD changes and the model is
regenerated, the numbers in this document change with it.

% =====================================================================
\section{System overview}

\archfig{system_architecture}{System architecture. The red band is the
hardware/software boundary: two abstract interfaces, \code{JointBackend} and
\code{CameraBackend}. Simulation and hardware differ only below it.}{1.0}

\subsection{Design targets}

\begin{center}\small
\begin{tabular}{@{}lll@{}}
\toprule
\textbf{Quantity} & \textbf{Target} & \textbf{Achieved} \\
\midrule
Total mass            & 10\,kg              & \RdMassTotal\,kg \\
Degrees of freedom    & 12 (HAA/HFE/KFE)    & 12 \\
Nominal stance height & ---                 & \RdStandHeight\,mm (\RdStandExtension\,\% of leg reach) \\
Static torque margin  & within continuous   & 62.5\,\% of \RdTorqueCont\,N$\cdot$m \\
Control rate          & ---                 & \RdControlRate\,Hz, sim and hardware alike \\
Actuator              & ROBSTRIDE02         & \RdTorqueCont/\RdTorquePeak\,N$\cdot$m cont/peak \\
\bottomrule
\end{tabular}
\end{center}

\subsection{Mass budget}

The CAD assembly was drawn around RMD-X8 actuators at 760.8\,g each. Substituting
ROBSTRIDE02 at \RdMassActuator\,g removes 4.57\,kg and lands the structure at
\RdMassStruct\,kg, leaving \RdMassPayload\,kg for electronics and
\RdMassCamera\,kg for the camera.

\begin{center}\small
\begin{tabular}{@{}lrl@{}}
\toprule
\textbf{Item} & \textbf{Mass (kg)} & \textbf{Note} \\
\midrule
CAD structure                       & 4.332 & body, brackets, legs, bearings \\
12 $\times$ ROBSTRIDE02             & \RdMassTwelve & \RdMassActuator\,g each \\
Battery, 12S LiPo                   & 0.700 & \RdVoltage\,V nominal \\
Onboard computer                    & 0.200 & \\
Power distribution and wiring       & 0.130 & \\
IMU                                 & 0.010 & \\
Camera (NUWA HP60C)                 & \RdMassCamera & separate link \\
\midrule
\textbf{Total}                      & \textbf{\RdMassTotal} & \\
\bottomrule
\end{tabular}
\end{center}

Per-link masses after aggregation: base \RdMassBase\,kg, hip \RdMassHip\,g,
thigh \RdMassThigh\,g, calf \RdMassCalf\,g. The whole-robot centre of mass sits
\RdComX\,mm from the body centre longitudinally.

\subsection{Coordinate systems and conventions}
\label{sec:conventions}

\paragraph{Frames.} REP-103 throughout: $x$ forward, $y$ left, $z$ up.
\code{base\_link} is at the body geometric centre. \code{odom} is the world
frame for the estimated base pose and the test-world markers. Camera optical
frames follow the ROS optical convention ($z$ forward, $x$ right, $y$ down),
related to their body frames by $\mathrm{rpy} = (-\pi/2,\, 0,\, -\pi/2)$.

\paragraph{Joint naming.} \code{\{FL,FR,RL,RR\}\_\{haa,hfe,kfe\}\_joint}.
HAA is hip abduction/adduction, HFE hip flexion/extension, KFE the knee.

\begin{decision}{Joint axes are not mirrored between left and right}
Every HAA axis is $+x$ and every HFE and KFE axis is $+y$, expressed in
\code{base\_link} orientation, for all four legs. The alternative --- mirroring
the right side so that ``positive'' means abduction everywhere --- was rejected
because it makes the URDF axes differ per leg and turns every sign question into
a per-leg question.

The consequence must be remembered: \textbf{positive HAA abducts the left legs
and adducts the right legs.} Sim-to-real then reduces to one sign per joint in
\code{robstride\_bus.yaml}, not a frame convention.
\end{decision}

\paragraph{Zero and sign.} All joints at zero puts the legs straight down, feet
at $(\pm 0.221,\ \pm 0.1445,\ -0.430)$\,m. Positive HFE swings the thigh
backward. Negative KFE folds the knee backward; KFE\,$=0$ is fully extended and
the upper limit, so the knee cannot hyperextend.

\begin{decision}{Named poses are specified as a base height, not joint angles}
Each foot hangs vertically below its own HFE axis. That one rule yields, at
once: identical joint angles on all four legs despite the non-mirrored axes; a
support polygon of \RdFootprintL\,$\times$\,\RdFootprintW\,mm centred within
\RdComX\,mm of the centre of mass; and the lowest peak joint torque of any
centred foot placement.

Placing the feet under the \emph{HAA} axes instead would need different front
and rear angles, because the HFE sits \RdHfeDx\,mm outboard of the HAA, and
would cost 73\,\% of the continuous rating on the rear knees instead of 62.5\,\%.
\end{decision}

% =====================================================================
\section{Mechanical model}

\subsection{CAD to URDF pipeline}

The Onshape export is a flat tree of 85 part-links with 12 revolute joints,
saved at an arbitrary mate pose. \code{tools/cad\_to\_model.py} turns it into a
canonical 13-body model:

\begin{enumerate}
\item forward-kinematics every part into the \code{root} frame;
\item cut the tree at the 12 revolute joints, giving 13 rigid bodies;
\item recover the pose-invariant leg geometry (link lengths, axis offsets);
\item solve the CAD mate angles against the canonical zero definition;
\item re-express every visual mesh and inertia tensor in canonical body frames;
\item substitute ROBSTRIDE02 mass properties for the CAD's RMD-X8 placeholders.
\end{enumerate}

The canonical chain reproduces the CAD joint positions to within
\textbf{9.9\,$\mu$m}, and the four legs agree with each other to under 2\,$\mu$m
--- the CAD is exactly symmetric, and the residual is mesh float precision.

\begin{decision}{The URDF is generated, never hand-edited}
\code{robot\_parameters.yaml} is the single source of truth. The Xacro, the
MuJoCo MJCF, the inverse kinematics and this document all read it. A
hand-maintained URDF drifts from the CAD, and a drift between the model the
physics uses and the model the IK uses presents as a control bug that is very
hard to attribute.
\end{decision}

\subsection{Kinematics}

\begin{center}\small
\begin{tabular}{@{}lrl@{}}
\toprule
\textbf{Parameter} & \textbf{Value} & \textbf{Meaning} \\
\midrule
HAA position       & $\pm$\RdHaaX, $\pm$\RdHaaY\,mm & from \code{base\_link} \\
HAA $\to$ HFE      & \RdHfeDx\,mm longitudinal, \RdHfeDr\,mm radial & \\
Thigh length $L_1$ & \RdThighLength\,mm & HFE to KFE, perpendicular \\
Thigh lateral      & \RdThighLat\,mm & along the HFE axis \\
Shank length $L_2$ & \RdShankLength\,mm & KFE to foot centre \\
Leg-plane offset   & \RdLateral\,mm & constant, $= $ radial $+$ lateral \\
Maximum reach      & \RdReach\,mm & $L_1 + L_2$ \\
Foot sphere        & \RdFootRadius\,mm radius & contact geometry \\
\bottomrule
\end{tabular}
\end{center}

Because rotation about $y$ leaves $y$ unchanged, the lateral offset collapses to
a constant. That is what makes the inverse kinematics closed-form: the abduction
angle is fixed by the $(y,z)$ projection alone, and what remains is a planar
two-link problem. Forward kinematics from the HAA joint frame:
\[
p = R_x(q_0)\bigl(d_1 + R_y(q_1)\bigl(d_2 + R_y(q_2)\, d_3\bigr)\bigr)
\]
with $d_1 = (s_x\,\text{hfe}_{dx},\, s_y\,\text{hfe}_{dr},\, 0)$,
$d_2 = (0,\, s_y\,\text{lat},\, -L_1)$, $d_3 = (0,0,-L_2)$.

The IK returns only the knee-backward branch ($q_2 \le 0$); the mirror solution
exists mathematically but the robot cannot reach it. Round-tripping FK through
IK over the full joint range gives a worst-case position error below
$10^{-9}$\,m, and the analytic Jacobian matches central differences to
$10^{-6}$.

\subsection{Knee drive topology}

The knee actuator is mounted \emph{proximally} on the thigh --- 48\,mm from the
HFE axis, 179\,mm from the knee --- and drives the joint through a 1:1 belt.

\begin{decision}{Proximal knee actuator, modelled as a simple revolute joint}
This is the standard low-distal-inertia arrangement. Its effect is measurable in
the generated model: the thigh's centre of mass sits only 27\,mm below the hip,
and the calf masses just \RdMassCalf\,g. Because the motor rotates \emph{with}
the thigh, the knee angle is not kinematically coupled to the hip angle, so the
URDF needs no mimic joint or transmission --- the belt ratio is a parameter that
is currently 1.0. Should a non-unity ratio be fitted, it enters in one place:
the joint-to-motor map.
\end{decision}

\subsection{Joint limits}

\begin{center}\small
\begin{tabular}{@{}lrrl@{}}
\toprule
\textbf{Joint} & \textbf{Lower (rad)} & \textbf{Upper (rad)} & \textbf{Rationale} \\
\midrule
HAA & \RdHaaLo & \RdHaaHi & $\pm$\RdHaaDeg$^\circ$ abduction envelope \\
HFE & \RdHfeLo & \RdHfeHi & thigh forward to folded \\
KFE & \RdKfeLo & \RdKfeHi & no hyperextension; folds backward only \\
\bottomrule
\end{tabular}
\end{center}

\begin{caveat}{The CAD's own joint limits are not used}
The Onshape export carries mate limits that are modelling artefacts --- for
example \code{dof\_fl0} spans $[-0.704, 0.866]$, asymmetric and different on
every leg. They describe how the assembly was posed, not what the mechanism can
do. The limits above are design values chosen from the mechanical envelope and
should be refined once self-collision checking runs against the real geometry.
\end{caveat}

% =====================================================================
\section{Actuator: ROBSTRIDE02}
\label{sec:actuator}

\subsection{Specification and derived quantities}

\begin{center}\small
\begin{tabular}{@{}llr@{}}
\toprule
\textbf{Quantity} & \textbf{Source} & \textbf{Value} \\
\midrule
Mass                    & spec    & \RdMassActuator\,g $\pm$ 3\,g \\
Envelope                & spec    & 78.5\,mm dia $\times$ 41.5\,mm \\
Gear ratio              & spec    & \RdGearRatio:1, machined steel \\
Rated / peak torque     & spec    & \RdTorqueRated\,/\,\RdTorquePeak\,N$\cdot$m \\
Rated speed             & spec    & 100\,rpm $=$ \RdSpeedRated\,rad/s \\
No-load speed           & spec    & 410\,rpm $=$ \RdSpeedNoLoad\,rad/s \\
Torque constant $K_t$   & spec    & \RdKt\,N$\cdot$m/A$_{\mathrm{rms}}$ \\
Line resistance         & spec    & 0.58\,$\Omega$ (phase \RdRphase\,$\Omega$) \\
Voltage                 & spec    & \RdVoltage\,V rated, \RdVoltMin--\RdVoltMax\,V \\
Encoders                & spec    & 2 $\times$ \RdEncoderBits-bit absolute magnetic \\
Output resolution       & derived & \RdEncoderRes\,rad \\
Rated phase current     & derived & \RdRatedArms\,A$_{\mathrm{rms}}$ \\
Copper loss at rating   & derived & \RdCopperLoss\,W \\
Reflected rotor inertia & estimate& \RdArmature\,kg$\cdot$m$^2$ \\
\bottomrule
\end{tabular}
\end{center}

\paragraph{$K_t$ is output-referred.} Dividing the rated torque by the torque
constant gives $\RdTorqueRated / \RdKt = 4.92\,\mathrm{A_{rms}}$, which matches
the rated phase current of 7\,A peak $/\sqrt{2} = \RdRatedArms\,\mathrm{A_{rms}}$.
The constant is therefore quoted \emph{after} the \RdGearRatio:1 reduction. This
matters: it makes the torque-to-current conversion used by the thermal observer
a single multiplication, with no gear ratio in it.

\paragraph{Reflected inertia dominates the distal link.} With
$J_{\mathrm{rotor}} \approx \RdRotorInertia\,\mathrm{kg\,m^2}$ estimated for a
28-pole outrunner, the inertia seen at the output is
$J_{\mathrm{rotor}} \times \RdGearRatio^2 = \RdArmature\,\mathrm{kg\,m^2}$ ---
larger than the calf's own inertia about the knee. Omitting it from the
simulation would make the legs roughly four times too easy to accelerate.

\subsection{CAN interface}

Extended 29-bit identifiers at \RdCanRate\,Mbit/s. The identifier carries
\texttt{[28:24] comm type | [23:8] data | [7:0] target}. The motion-control
frame (type 1) is the MIT-style impedance command: target angle, velocity, $k_p$
and $k_d$ occupy the eight payload bytes, and the feed-forward torque rides in
the identifier's data field because the payload is full.

\begin{decision}{The command interface is the actuator's native impedance frame}
\code{JointCommand} carries $(\text{position},\ \text{velocity},\ \text{effort},\
k_p,\ k_d)$ --- exactly one CAN frame. A narrower interface, such as position
only, would force the host to close a loop the motor already closes in firmware
at tens of kHz. The host sets \emph{setpoints}; it is not the servo.
\end{decision}

\subsection{Bus budget}

An extended-ID frame with eight data bytes is 131 bits nominal and up to 159
with worst-case stuffing; \RdBitsPerFrame\ bits is the design figure. Each joint
needs one command and one feedback frame per cycle.

\begin{center}\small
\begin{tabular}{@{}rrrrl@{}}
\toprule
\textbf{Buses} & \textbf{Motors/bus} & \textbf{Rate} & \textbf{Load} & \textbf{Verdict} \\
\midrule
1 & 12 & 250\,Hz & 90\,\% & too tight, no margin \\
1 & 12 & 400\,Hz & 144\,\% & impossible \\
2 & 6  & \RdControlRate\,Hz & \RdBusLoad\,\% & \textbf{selected} \\
2 & 6  & 500\,Hz & 90\,\% & rejected, no margin \\
3 & 4  & 500\,Hz & 60\,\% & upgrade path \\
\bottomrule
\end{tabular}
\end{center}

\begin{decision}{Two buses at \RdControlRate\,Hz, split front/rear}
\RdControlRate\,Hz is not a compromise on control quality: because the RS02
closes its position and velocity loop in firmware, the host rate is a
\emph{setpoint update} rate. The same rate is used in simulation so that gains
tuned in MuJoCo transfer without change.

The split is front/rear rather than left/right so that each bus carries a
complete leg \emph{pair}: a single bus failure leaves the robot in a
recognisable posture rather than a diagonal collapse.
\code{RobStride02Backend.check\_bus\_budget()} refuses any configuration above
80\,\% at start-up.
\end{decision}

% =====================================================================
\section{Software architecture}

\archfig{software_architecture}{ROS 2 package graph. Arrows point from
dependant to dependency. No package above the hardware layer depends on a
simulation package.}{0.96}

\subsection{Packages}

\begin{center}\small
\begin{tabular}{@{}lll@{}}
\toprule
\textbf{Package} & \textbf{Build} & \textbf{Responsibility} \\
\midrule
\code{robodog\_msgs}        & ament\_cmake  & message and service contract \\
\code{robodog\_description} & ament\_cmake  & URDF/Xacro, meshes, RViz, actuator data \\
\code{robodog\_hardware}    & ament\_python & \code{JointBackend} + RS02 CAN codec \\
\code{robodog\_control}     & ament\_python & control loop, safety, gait, IK, estimator \\
\code{robodog\_sim}         & ament\_python & MJCF generator, test worlds, markers \\
\code{robodog\_perception}  & ament\_python & \code{CameraBackend}, depth, point cloud \\
\code{robodog\_web}         & ament\_python & GUI server and protocol \\
\code{robodog\_bringup}     & ament\_cmake  & launch composition only \\
\bottomrule
\end{tabular}
\end{center}

\code{robodog\_msgs} depends on nothing but \code{std\_msgs} and
\code{geometry\_msgs}: it is the contract between layers and must not acquire an
implementation dependency.

\subsection{The hardware boundary}
\label{sec:boundary}

\archfig{hardware_boundary}{Everything above the boundary is unchanged by the
swap to hardware. The per-joint direction and offset map is the only place where
motor coordinates and canonical joint coordinates differ.}{1.0}

\code{JointBackend} is a six-method abstract class:

\begin{lstlisting}[language=Python]
class JointBackend(ABC):
    name: str; is_simulation: bool
    def configure(self) -> None          # open devices / load model
    def enable(self, mask) -> None       # energise
    def disable(self, mask) -> None      # de-energise, back-driveable
    def read(self) -> JointState         # pos, vel, effort, temperature, faults
    def write(self, cmd: JointCommand)   # already safety-clamped
    def step(self, dt) -> None           # advance a simulator; no-op on hardware
    def base_state(self) -> BaseState | None   # None on real hardware
\end{lstlisting}

\begin{decision}{\code{base\_state()} returns \code{None} on the real robot}
A simulator knows the floating-base state exactly; the real robot has no sensor
for absolute base pose. Rather than inventing one, the contract makes the
absence explicit and forces every consumer to handle it. The state estimator
fills the gap from the IMU and leg kinematics, and marks its output
\code{ground\_truth = False}. Returning a plausible-looking zero would hide the
difference that matters most when a controller tuned in simulation is first run
on hardware.
\end{decision}

Three implementations exist. \code{KinematicBackend} models each joint as an
independent second-order system with the reflected rotor inertia --- no physics
engine, no contact, no gravity --- so the entire stack above the boundary is
runnable and testable with nothing installed. \code{MujocoBackend} adds
rigid-body dynamics and contact. \code{RobStride02Backend} drives the real
motors.

\subsection{Interfaces}

\archfig{node_topic_graph}{Node and topic graph. Identical in simulation and on
hardware; only the implementations behind the two backends differ.}{1.0}

\begin{center}\small
\begin{tabular}{@{}llrl@{}}
\toprule
\textbf{Topic} & \textbf{Type} & \textbf{Rate} & \textbf{Purpose} \\
\midrule
\topicname{/joint\_states}              & JointState        & 100 & TF, RViz, standard tooling \\
\topicname{/robodog/robot\_state}       & RobotState        & 50  & aggregate telemetry \\
\topicname{/robodog/safety\_status}     & SafetyStatus      & 50  & limits, faults, loop health \\
\topicname{/robodog/imu}                & Imu               & 50  & attitude and rates \\
\topicname{/robodog/joint\_command}     & JointCommandArray & --- & direct joint control \\
\topicname{/robodog/gait\_command}      & GaitCommand       & --- & gait selection \\
\topicname{/cmd\_vel}                   & Twist             & --- & body velocity \\
\topicname{/robodog/world\_markers}     & MarkerArray       & latched & test environment \\
\topicname{camera/color/image\_raw}     & Image (rgb8)      & 15  & colour \\
\topicname{camera/depth/image\_rect\_raw}& Image (16UC1)    & 15  & depth, millimetres \\
\topicname{camera/depth/points}         & PointCloud2       & 10  & range-filtered cloud \\
\bottomrule
\end{tabular}
\end{center}

Services: \code{set\_named\_pose}, \code{set\_gait}, \code{set\_control\_mode},
\code{enable\_joints}, \code{emergency\_stop}.

\begin{decision}{One aggregate telemetry message}
\code{RobotState} carries all twelve joints, the four feet, the base estimate,
controller state, safety state and simulation state in one message. A consumer
therefore never has to time-align several topics to render a consistent picture.
\code{/joint\_states} is published separately and at a different rate because
the wider ROS ecosystem expects it.
\end{decision}

% =====================================================================
\section{Control architecture}

\archfig{control_architecture}{Control layers. Rates fall by roughly an order of
magnitude per layer, and the innermost loop is not on the host at all.}{0.98}

\subsection{The loop}

One thread, \RdControlRate\,Hz, in \code{robodog\_control\_node}:

\begin{center}\small
\code{read()} $\to$ estimate $\to$ active controller $\to$ safety $\to$ \code{write()} $\to$ \code{step()}
\end{center}

ROS callbacks run on a separate executor and only ever set fields the loop
reads; nothing in a callback touches the backend. Publishing is decimated ---
\code{RobotState} at 50\,Hz, \code{/joint\_states} at 100\,Hz --- because
publishing telemetry at \RdControlRate\,Hz would cost more than the control it
reports on.

\archfig{data_flow}{One \RdControlRate\,Hz cycle. The safety monitor is
unconditional.}{1.0}

\begin{decision}{Controllers are selected, not composed}
Exactly one of \{\code{idle}, \code{joint}, \code{pose}, \code{gait},
\code{joint\_test}\} drives the joints at any moment. Making that exclusive by
construction is worth more than the flexibility of a controller chain, because
two things writing to the same actuator is a failure mode with no good recovery.
\end{decision}

\subsection{Impedance control}

The actuator evaluates
\[
\tau = k_p\,(q^{*} - q) + k_d\,(\dot q^{*} - \dot q) + \tau_{\mathrm{ff}}
\]
in firmware. Effective joint inertia is the reflected rotor inertia
(\RdArmature) plus roughly $2\times10^{-3}$ of link, so $J \approx
7\times10^{-3}\,\mathrm{kg\,m^2}$; $k_p = 120$ then gives
$\omega_n = \sqrt{k_p/J} \approx 131$\,rad/s (21\,Hz) and a critical damping of
$k_d = 2\sqrt{k_p J} = 1.83$.

\begin{center}\small
\begin{tabular}{@{}lrrl@{}}
\toprule
\textbf{Context} & $k_p$ & $k_d$ & \textbf{Rationale} \\
\midrule
Named pose, joint test & 80  & 2.0 & moderate, unloaded \\
Stance leg             & 120 & 2.5 & stiff enough to hold the body up \\
Swing leg              & 45  & 1.2 & compliant: an unexpected contact during
                                     swing must not lever the robot over \\
\bottomrule
\end{tabular}
\end{center}

\subsection{Gravity feed-forward and the force convention}

Each stance leg is fed the torque that holds its share of the body weight. The
sign convention is stated once and used everywhere:

\begin{center}
$f$ is the force the \emph{environment} applies to the foot (the ground
reaction, $+z$ in stance), so \quad
$\tau_{\mathrm{actuator}} = -J^{\mathsf{T}} f$, \quad
$f = -(J^{\mathsf{T}})^{-1}\tau_{\mathrm{actuator}}$.
\end{center}

\begin{caveat}{This sign is not cosmetic}
An inverted feed-forward makes a stance leg push the body \emph{down}. Measured
in MuJoCo: with correct feed-forward the body sags 0.31\,mm below the commanded
\RdStandHeight\,mm; without feed-forward it sags 5.34\,mm; with the sign
inverted it sags 11.95\,mm. Magnitude-only tests pass in all three cases, which
is why \code{test\_kinematics.py} asserts the \emph{direction} of the knee
torque explicitly.
\end{caveat}

\subsection{Gait generation}

A single normalised cycle phase advances at the step frequency. Each leg has a
fixed offset and is in stance for \code{duty\_factor} of its period.

\begin{center}\small
\begin{tabular}{@{}llrrl@{}}
\toprule
\textbf{Gait} & \textbf{Offsets} & \textbf{Duty} & \textbf{Freq (Hz)} & \textbf{Character} \\
\midrule
stand & ---                     & 1.00 & 0.0 & all feet planted \\
walk  & 0, 0.50, 0.75, 0.25     & 0.75 & 1.2 & statically stable, one foot up \\
trot  & 0, 0.50, 0.50, 0        & 0.50 & 2.2 & diagonal pairs; default travel \\
pace  & 0, 0.50, 0, 0.50        & 0.50 & 2.0 & lateral pairs; rolls this narrow body \\
bound & 0, 0, 0.50, 0.50        & 0.45 & 2.6 & front then rear; most demanding \\
\bottomrule
\end{tabular}
\end{center}

\subsubsection{What the generator actually does}

The gait generator is a \emph{foot trajectory planner}. It does not compute
torques and it does not know what a motor is. Once per control cycle it is
given the elapsed time and the body's measured state, and it returns, for each
of the four legs, where that foot should be, how fast it should be going, and
how hard it should push. Everything else --- inverse kinematics, the impedance
law, the actuator --- happens downstream.

Concretely, \code{GaitGenerator.update(dt, feedback)} returns twelve joint
positions, twelve joint velocities, twelve feed-forward torques, and a
four-element contact flag. The joint values come from running each foot
position through the closed-form leg IK; the torques come from the balance
layer. A gait is therefore defined entirely by \emph{geometry and timing}, and
adding one is a block in \code{gaits.yaml} plus a phase offset --- no new code.

\archfig{gait_pipeline}{One control cycle of the gait generator. The
generator plans feet, not torques: everything it returns is a position, a
velocity or a force, and the impedance law runs two layers below. Continuity
is the invariant --- stance resumes from the last commanded position and swing
starts from where stance left the foot, so the command cannot step at either
transition.}{1.0}

\subsubsection{The phase model}

One normalised phase $\phi \in [0,1)$ advances at the step frequency and wraps.
Each leg $i$ reads the same clock through its own offset:
\[
  \phi_i = (\phi + \text{offset}_i) \bmod 1,
  \qquad
  \text{leg } i \text{ is in \textsc{stance} iff } \phi_i < D
\]
where $D$ is the duty factor. That single line is what distinguishes a walk
from a trot: the offsets decide which feet are down together, and $D$ decides
for how long. A duty of 0.75 with the walk offsets always leaves three feet
down, so the robot is statically stable; a duty of 0.5 with the trot offsets
leaves two, on a diagonal, so it is not.

\subsubsection{Stance: the foot is a fixed point in the world}

A planted foot does not move. It follows that in the \emph{body} frame it must
translate backwards at exactly the body's velocity --- if the body advances
10\,mm, the foot must appear 10\,mm further back, or the command is asking the
leg to drag the robot through the contact.

\[
  \mathbf{p}^{\text{stance}}_i \mathrel{-}= \mathbf{v}\,\delta t,
  \qquad
  \mathbf{v} = \mathbf{v}_{\text{meas}}
    + g_{\text{sweep}}\,(\mathbf{v}^{*} - \mathbf{v}_{\text{meas}})
\]

The sweep gain $g_{\text{sweep}}$ chooses which velocity to believe, and the
default of 1.0 --- the \emph{commanded} velocity --- is the least obvious
decision in the whole generator. It looks like it must scuff. In steady state
it does not: once the body is travelling at the commanded speed, the target
sweeps at exactly the rate the planted foot is already moving through the body
frame, and the relative motion is zero. When the body is \emph{not} at speed,
the target runs ahead of the foot and the leg impedance turns that gap into a
tangential force.

That is the point. Sweeping at the commanded velocity is proportional velocity
feedback routed through the leg, and it is the dominant thing regulating speed.
Sweeping at the measured velocity instead is self-consistent at any speed,
including the wrong one: measured at $g_{\text{sweep}} = 0$, a walk commanded
to 0.15\,m/s free-ran to 0.272\,m/s and then fell over.

\subsubsection{Swing: interpolate between two latched endpoints}

Swing carries the foot from where it left the ground to where it should next
land. Both endpoints are explicit positions:

\begin{itemize}
  \item $\mathbf{p}_0$ is latched at lift-off --- wherever stance actually left
    the foot, not where it was supposed to be;
  \item $\mathbf{p}_1$ is the landing point, tracked continuously through the
    swing so a foot in the air still reacts to the body accelerating.
\end{itemize}

\[
  \mathbf{p}^{\text{swing}}(u)
    = \mathbf{p}_0 + (\mathbf{p}_1 - \mathbf{p}_0)\,\sigma(u),
  \qquad
  \sigma(u) = u - \frac{\sin 2\pi u}{2\pi}
\]

$\sigma'(u) = 1 - \cos 2\pi u$ vanishes at both ends, so the foot leaves and
arrives with no horizontal velocity relative to the ground --- it is placed,
not dragged. The vertical profile $\sin^2(\pi u)$ is likewise zero-slope at
both ends, with a small \code{touchdown\_depth} subtracted late in swing so the
foot seeks the floor rather than hovering over it, and decayed again over the
first sixth of stance so that seeking does not itself become a step.

Interpolating between two latched endpoints is what makes the command
continuous, and it is worth being explicit about why that is not a detail.
Section~\ref{sec:locomotion-status} records what happened when it was not.

\subsubsection{Foot placement: where velocity is really decided}

Attitude control holds a robot up; it cannot make it go anywhere. A legged
robot regulates velocity by \emph{where it puts its feet}. The landing point
uses the Raibert heuristic:
\[
  \mathbf{p}_1 = \mathbf{p}_{\text{nom}}
    + \underbrace{\mathbf{v}_{\text{meas}}\,\frac{T_{\text{stance}}}{2}}_{\text{neutral point}}
    + \underbrace{k_v\,(\mathbf{v}_{\text{meas}} - \mathbf{v}^{*})}_{\text{correction}}
\]
The first term is the landing point that \emph{sustains} the current velocity:
land half a stance-sweep ahead, sweep back a full stance-sweep, end half behind,
and the stride is symmetric about the hip. The second term breaks that symmetry
deliberately --- landing further ahead than neutral decelerates, further behind
accelerates. The offset is clamped to \code{max\_placement\_m} so a large
velocity error cannot ask for a footfall outside the workspace.

\subsubsection{Three things regulate speed, and they are not interchangeable}

\begin{center}\small
\begin{tabular}{@{}llp{0.42\textwidth}@{}}
\toprule
\textbf{Mechanism} & \textbf{Acts through} & \textbf{Role} \\
\midrule
Stance sweep     & leg impedance & Dominant. Turns velocity error into
                                   tangential force every cycle. \\
Tangential wrench & feed-forward torque & Explicit force request,
                                   $k_{p,v}(\mathbf{v}^{*}-\mathbf{v})$,
                                   friction-cone limited. \\
Foot placement   & next footfall & Sets the velocity the gait is stable
                                   \emph{at}; acts once per step, not
                                   continuously. \\
\bottomrule
\end{tabular}
\end{center}

Removing any one degrades tracking; removing the first breaks it outright.

\subsubsection{Force distribution}

Positions alone do not hold a robot up. The contact flag from the phase model
tells the balance layer which feet may push, and it solves for the ground
reaction at each of them that best reproduces the desired body wrench
(Section~\ref{sec:balance}). Those forces become feed-forward joint torques
through $\boldsymbol{\tau} = -J^{\top}\mathbf{f}$. The gait generator supplies
the geometry and the contact schedule; the balance layer supplies the forces;
neither one alone walks.

\subsubsection{Per-cycle summary}

\begin{enumerate}\itemsep2pt
  \item Advance $\phi$; for each leg compute $\phi_i$ and its stance/swing state.
  \item Compute the Raibert landing offset from measured and commanded velocity.
  \item \textsc{Stance}: on touchdown, resume from the last commanded position;
        retract at the sweep velocity; decay the touchdown depth.
  \item \textsc{Swing}: on lift-off, latch $\mathbf{p}_0$; interpolate
        $\mathbf{p}_0 \to \mathbf{p}_1$ with the zero-endpoint-velocity profile
        and the $\sin^2$ lift.
  \item Run each foot position and velocity through the leg IK and Jacobian to
        get $q$ and $\dot q$.
  \item Ask the balance layer for the body wrench, distribute it over the
        stance feet, convert to feed-forward torque.
  \item Return $q$, $\dot q$, $\tau_{\text{ff}}$, contact.
\end{enumerate}

Stride length is never configured, and there is no stride parameter to
configure. Swing interpolates between where the foot left the ground and where
placement says it should land, so the stride is whatever those two points
imply. The workspace cap is \code{max\_placement\_m}, which bounds the landing
offset directly --- 0.12\,m against a 0.43\,m leg, comfortably inside reach.

\subsection{Balance}
\label{sec:balance}

A foot-trajectory generator alone does not walk. With the open-loop gait and no
balance term, a 0.3\,m/s trot in MuJoCo reached 25$^\circ$ of body tilt,
saturated the actuators at the \RdTorquePeak\,N$\cdot$m peak and travelled
28\,mm in six seconds. The feet were following their commanded paths; the body
was not being held up. \code{balance.py} adds the two things that were missing.

\paragraph{Force distribution.} A PD controller on height, roll, pitch, yaw rate
and horizontal velocity produces a desired body wrench, and the stance foot
forces are solved to reproduce it:
\[
\sum_i f_i = F, \qquad \sum_i r_i \times f_i = M
\]
Six equations in $3 n_{\mathrm{stance}}$ unknowns --- exactly determined for a
trot, under-determined for a walk or a full stance --- solved by least squares,
then projected into what a contact can deliver: $f_z \ge 0$ and
$\lVert f_{xy} \rVert \le \mu f_z$.

\begin{decision}{A stance leg is force controlled, not position controlled}
This was measured the hard way. With \code{stance\_kp} at 120 the stance legs
showed 0.17\,rad of tracking error and 20\,N$\cdot$m of pure disagreement with
reality: the commanded foot position was advancing at the \emph{commanded} body
velocity while the foot was physically planted, so the leg was dragging the
robot through its own contact. Two changes follow. The stance foot now
translates at a blend of the \emph{measured} and commanded velocity
(\code{stance\_sweep\_gain}), and the placement offset is latched at touchdown
so the command cannot move under a planted foot. Stiffness is reduced to 90:
raising it turns the leg into a spring that anchors the body to its own foot
and cancels the propulsion force.
\end{decision}

\paragraph{Foot placement.} Attitude control cannot regulate velocity --- a
legged robot steers by where it puts its feet. The Raibert heuristic places the
landing point at $v\,T_{\mathrm{stance}}/2 + k_v (v - v^{*})$.

\subsection{Measured locomotion performance}
\label{sec:locomotion-status}

Six seconds of each gait from the standing keyframe, commanded in a straight
line, with the balance layer active. Drift is the largest lateral excursion.

\begin{center}\small
\begin{tabular}{@{}lrrrrrl@{}}
\toprule
\textbf{Gait} & \textbf{$v^{*}_x$} & \textbf{actual} & \textbf{error} &
\textbf{tilt} & \textbf{drift} & \textbf{Verdict} \\
 & (m/s) & (m/s) & (m/s) & (deg) & (m) & \\
\midrule
stand & 0.00 & $-0.001$ & $-0.001$ &  0.1 & 0.000 & solid \\
walk  & 0.15 & $+0.155$ & $+0.005$ &  7.4 & 0.024 & tracks \\
walk  & 0.30 & $+0.300$ & $+0.000$ & 10.9 & 0.040 & tracks \\
trot  & 0.30 & $+0.297$ & $-0.003$ &  5.1 & 0.040 & tracks; default \\
trot  & 0.50 & $+0.493$ & $-0.007$ &  7.2 & 0.054 & tracks \\
pace  & 0.30 & $+0.146$ & $-0.154$ & 16.6 & 0.066 & rolls; comparison only \\
bound & 0.30 & $+0.193$ & $-0.107$ & 18.4 & 0.001 & probe gait only \\
\bottomrule
\end{tabular}
\end{center}

Walk and trot --- the two gaits intended for travel --- track commanded
velocity to within 7\,mm/s. Pace and bound do not, which is consistent with
what they are for: pace rolls a body only 120\,mm wide between hip
pivots and is kept for comparison, and bound exists to probe the peak-torque
and $I^2t$ limits in simulation before hardware ever sees them.

\subsubsection{How it was broken, and what that cost}

For most of this project the robot did not travel. It stepped in place and
drifted backwards at roughly 0.04\,m/s, while every travelling gait drew 141
to 201\,\% of the continuous torque rating. Both symptoms had the same two
causes, and neither was in the balance layer or the gains, which is where the
effort went first.

\paragraph{A 41\,mm step at every lift-off.} Stance retracts the foot at the
sweep velocity, but the swing arc was rebuilt from scratch each cycle as
$\mathbf{p}_{\text{nom}} + \text{stride}\,(\sigma - \tfrac12)$ --- which
assumes stance had swept the foot back by the full \emph{commanded} stride.
Whenever the robot was not already at the commanded speed the two disagreed,
and the difference was commanded in a single 2.5\,ms tick at every lift-off:
41\,mm from a standstill at 0.30\,m/s. That step drove 0.33\,rad of tracking
error, which saturated the hips at the \RdTorquePeak\,N$\cdot$m peak and so
destroyed the force distribution the balance layer had just computed. The robot
could not accelerate, so the measured velocity stayed at zero, so the step
never shrank. The gait fought itself into a steady state.

\paragraph{Two jammed hips.} The thigh capsule and the HAA actuator cylinder
are both conservative bounding volumes around parts that clear each other in
the CAD, and MuJoCo does not auto-filter the pair because the thigh is the
base's \emph{grandchild} (base $\to$ hip $\to$ thigh), not its child. The
front hips therefore jammed 13\,mrad above the standing pose and would not
move under 16.6\,N$\cdot$m, while the rear hips swung 1.59\,rad on 3\,N$\cdot$m.
The front legs could only lift with the knee, and the hip absorbed the entire
impedance error as heat. The robot was, in the most literal sense, walking with
two of its twelve joints welded shut.

\begin{caveat}{A careful measurement of a broken system is still a measurement of a broken system}
The stance-height study in \code{gaits.yaml} was repeated at four heights, was
internally consistent, and was wrong --- it concluded that the balance layer
stopped helping above 340\,mm, which was an artifact of the jam. So was the
finding that lower stance impedance reduced torque: with the hips free, the
original \code{stance\_kp} of 90 turned out to be the best value tested. The
tuning was never the problem, and no amount of it would have helped.

What eventually found both bugs was not more tuning but per-joint
instrumentation: printing commanded against measured position for one leg
through one gait cycle showed a 40\,mm jump in the target at lift-off, and a
hip angle that did not change by 0.001\,rad across four seconds under peak
torque. A joint that will not move under 16\,N$\cdot$m is not a control
problem.
\end{caveat}

Both are now covered by tests that fail if they return:
\code{test\_the\_foot\_command\_never\_jumps} asserts the commanded foot
position never moves more than 5\,mm in one tick, at four different measured
velocities, and
\code{test\_every\_hip\_can\_swing\_through\_its\_commanded\_range} drives each
hip alone and requires all four to move, and to move alike.

\subsubsection{What the balance layer is for, now}

Fixing the trajectory changed the answer. The balance layer used to be what
kept a 0.3\,m/s trot from reaching 25$^\circ$ of tilt; with the swing arc
continuous and the hips free, the open-loop gait is well behaved on its own and
the two are within a few tenths of a degree --- 5.1$^\circ$ with balance
against 4.9$^\circ$ without. Tilt is no longer where the layer earns its place.

Course holding is. Over six seconds at 0.5\,m/s the lateral drift is 0.054\,m
with balance and 0.299\,m without, because the Raibert term is the only thing
steering the feet back under the body, and velocity tracking degrades from
$-0.007$ to $-0.018$\,m/s. A robot that wanders 0.3\,m sideways in six seconds
does not get through a doorway. The regression test asserts those two things
rather than tilt.

\subsection{Thermal load}
\label{sec:thermal-load}

Peak torque is the wrong question to ask of a motor. What decides whether a
winding survives is the RMS torque, because copper loss goes as current
squared: a \RdTorquePeak\,N$\cdot$m spike lasting 20\,ms is harmless, and
10\,N$\cdot$m sustained is not. Measured over six seconds of each gait, worst
joint, with the start-up transient discarded:

\begin{center}\small
\begin{tabular}{@{}lrrrrrl@{}}
\toprule
\textbf{Gait} & \textbf{peak} & \textbf{p99} & \textbf{RMS} &
\textbf{RMS / cont.} & \textbf{$T_\infty$} & \textbf{Verdict} \\
 & (N$\cdot$m) & (N$\cdot$m) & (N$\cdot$m) & & ($^\circ$C) & \\
\midrule
stand         &  3.9 &  3.9 &  3.9 &  65\,\% &  45 & sustainable \\
walk 0.15     & 17.0 &  9.6 &  5.5 &  91\,\% &  70 & sustainable \\
trot 0.30     & 14.3 &  8.7 &  5.4 &  90\,\% &  69 & sustainable \\
trot 0.50     & 14.5 &  9.3 &  5.5 &  92\,\% &  71 & sustainable \\
bound 0.30    & 17.0 & 16.7 &  7.9 & 131\,\% & 123 & over limit \\
\bottomrule
\end{tabular}
\end{center}

These supersede an earlier table in which every travelling gait sat between 141
and 201\,\% of continuous, implying winding temperatures from 139 to
260\,$^\circ$C. That table was correct about a robot with two welded hips and a
discontinuous swing arc.

The interpretation was correct then and remains the useful one: the actuators
were never undersized. Standing costs 3.9\,N$\cdot$m RMS, and steady walking at
0.5\,m/s needs about 49\,W of net mechanical work against 754\,W rated --- some
7\,\% --- because almost all the torque a quadruped spends is holding itself up,
which does no work at all. A travelling gait that costs three times the RMS of
standing still is not doing three times the work; it is wasting the difference.
With the waste removed, trot sits at 90\,\% of continuous and 69\,$^\circ$C,
inside the \RdTempWarn\,$^\circ$C warning threshold.

Bound remains over limit, which is what it is for: it is the gait used to probe
the peak-torque and $I^2t$ path in simulation, and it should not be commanded on
hardware. Pace is not listed because it does not track velocity well enough for
the number to mean anything.

\begin{caveat}{Margin, not comfort}
90\,\% of continuous is passing, not roomy. Trot buys its margin partly by
standing 340\,mm tall rather than 320 --- a straighter leg has a shorter moment
arm --- and that costs 20\,mm of the leg travel otherwise available for terrain
and attitude correction. Walk stays at 320\,mm, where it is 91\,\%,
because it is the gait for stairs and narrow passages and wants the travel.
The per-gait stance heights in \code{gaits.yaml} carry the measured trade-off.

None of this has been validated against hardware. The thermal model is a
first-order estimate from datasheet resistance and an assumed thermal
resistance; the friction and damping are estimates. Treat 90\,\% as
``plausibly inside the rating'', not as a measurement of a real motor.
\end{caveat}

\subsection{State estimation}

\begin{decision}{A complementary filter and leg odometry, not an EKF}
Attitude comes from blending gyro integration with the gravity direction, with
the gravity term weighted by how closely the measured acceleration magnitude
resembles $g$ --- during a hard push-off it does not. Height and velocity come
from treating each foot in contact as a temporary fixed point.

An EKF was rejected \emph{for now} because until the robot exists its process
and measurement covariances would be invented numbers. This is honest,
debuggable and sufficient to stand, walk and visualise. Yaw is not observable
from an IMU without a magnetometer and is flagged as drifting.
\end{decision}

% =====================================================================
\section{Safety}
\label{sec:safety}

\subsection{Two-tier limits}

\begin{center}\small
\begin{tabular}{@{}llll@{}}
\toprule
\textbf{Quantity} & \textbf{Hard (URDF)} & \textbf{Operational} & \textbf{Enforced by} \\
\midrule
Torque   & \RdTorquePeak\,N$\cdot$m & \RdTorqueCont\,N$\cdot$m continuous & $I^2t$ integrator \\
Velocity & \RdSpeedNoLoad\,rad/s & \RdSpeedOp\,rad/s & clamp \\
Position & joint end stops & end stops inset 0.05\,rad & clamp + rate limit \\
Winding temperature & --- & \RdTempWarn/\RdTempDerate/\RdTempFault\,$^\circ$C & derate, then latch \\
\bottomrule
\end{tabular}
\end{center}

The URDF carries what the hardware can physically deliver; the operational tier
is what the robot is allowed to use.

\subsection{Design rules}

\begin{enumerate}
\item \textbf{It runs in simulation too.} The monitor executes on every cycle
      regardless of backend. A safety layer that only activates on hardware has
      never been tested by the time it matters.
\item \textbf{It is a pure function.} \code{(measurement, request, state)}
      $\to$ \code{(allowed command, status)}. No ROS, no I/O, no threads --- which
      is what makes the whole envelope testable. \code{test\_safety.py} is the
      specification.
\item \textbf{It clamps rather than rejects.} Dropping a cycle mid-stride is
      more dangerous than saturating one.
\item \textbf{Protective faults latch.} Clearing requires the condition to have
      gone away \emph{and} an explicit operator action.
\end{enumerate}

\subsection{Torque limiting under impedance control}

\begin{decision}{Scale the whole $(k_p, k_d, \tau_{\mathrm{ff}})$ triple}
The motor computes $k_p e + k_d \dot e + \tau_{\mathrm{ff}}$ itself, so clamping
only the \code{effort} field would not limit the torque at all. The monitor
predicts the resulting torque and, if it exceeds the budget, scales all three
terms by the same factor --- limiting the magnitude while preserving the
commanded impedance direction.
\end{decision}

The budget is the continuous rating normally, rising to peak while $I^2t$ credit
remains. The integrator charges as $(\tau/\tau_{\mathrm{cont}})^2 - 1$ above the
continuous rating, so twice the continuous torque exhausts the budget four times
faster, and discharges when below.

\subsection{Thermal model}

\[
P_{\mathrm{cu}} = 3 I_{\mathrm{rms}}^2 R_{\phi},\quad I_{\mathrm{rms}} = \tau/K_t,
\qquad
C_{\mathrm{th}}\,\dot T = P_{\mathrm{cu}} - \frac{T - T_{\mathrm{amb}}}{R_{\mathrm{th}}}
\]

with $R_{\mathrm{th}} = \RdRthermal$\,K/W and $C_{\mathrm{th}} = \RdCthermal$\,J/K,
giving a time constant of \RdTauThermal\,minutes. It runs in two roles: in
simulation it synthesises a plausible temperature, because MuJoCo has no concept
of copper loss and the safety architecture must be exercised; on the real robot
it runs as an \emph{observer} alongside the reported temperature, and the
monitor takes the maximum of the two. The motor's own sensor sits on the driver
board and lags the winding by tens of seconds --- the observer is what catches a
stalled leg pushing against an obstacle.

Between \RdTempDerate\ and \RdTempFault\,$^\circ$C the torque budget is scaled
down linearly rather than cut at a threshold, because cutting output abruptly
would drop the robot.

\subsection{Watchdog and emergency stop}

\begin{decision}{A stale command holds position; it does not go limp}
Losing torque mid-stance drops a 10\,kg robot. The watchdog substitutes an
impedance command that holds the measured position at moderate gain. If a joint
is measured outside its soft range, the hold target is the limit, so the
watchdog recovers the envelope gently rather than freezing outside it.
\end{decision}

A communication timeout on any joint latches an e-stop: a silently stale joint
is the failure mode most likely to break a leg.

\archfig{state_machine}{Robot states. FAULT and E-STOP are reachable from
everywhere; leaving E-STOP requires the cause to have cleared.}{0.86}

% =====================================================================
\section{Simulation architecture}

\archfig{simulation_architecture}{The MJCF and the URDF are generated from the
same parameters, so the physics and the kinematics cannot drift apart.}{1.0}

\subsection{Fidelity choices}

These are not convenience settings; each is here because getting it wrong yields
a controller that works in simulation and fails on the robot.

\begin{enumerate}
\item \textbf{The impedance law is evaluated at the physics rate}, not the
      control rate. The real RS02 closes its loop in firmware at tens of kHz;
      holding a \RdControlRate\,Hz zero-order hold on $k_p(q^*-q)$ would make the
      simulated joints noticeably softer than the hardware and would tune gains
      that are wrong on the robot. Only the \emph{setpoint} is held.
\item \textbf{\code{armature} carries the reflected rotor inertia}
      (\RdArmature\,kg$\cdot$m$^2$), which is larger than the calf's own inertia
      about the knee. Omitting it is not a small error.
\item \textbf{Command latency is injected} --- two physics steps, about 1\,ms,
      the order of a 1\,Mbit/s bus at \RdControlRate\,Hz. Controllers tuned
      against a zero-latency plant are optimistic in the wrong direction.
\item \textbf{Capsules, not cylinders}, for the leg segments. URDF has no
      capsule primitive; MuJoCo does, and capsule contact is far better
      conditioned than the flat-ended cylinder a URDF must approximate it with.
\item \textbf{Timestep 0.5\,ms} (2\,kHz), exactly five substeps per control
      cycle. A non-integer ratio would make the effective command latency jitter.
\end{enumerate}

Not modelled: gearbox backlash, belt compliance in the knee drive, current-loop
bandwidth, and CAN jitter.

\subsection{Validation against the model}

Standing in MuJoCo with gravity feed-forward, measured:

\begin{center}\small
\begin{tabular}{@{}lrr@{}}
\toprule
\textbf{Quantity} & \textbf{Predicted} & \textbf{Measured in MuJoCo} \\
\midrule
Total mass                & \RdMassTotal\,kg & \RdMassTotal\,kg \\
Ground reaction, all feet & 98.10\,N & 98.10\,N \\
Front load share          & 47.2\,\% & 47.0\,\% \\
Body height               & \RdStandHeight\,mm & 320.31\,mm \\
Peak joint torque         & 62.5\,\% of continuous & 64.6\,\% \\
\bottomrule
\end{tabular}
\end{center}

The front/rear split follows from the whole-robot centre of mass sitting
12.3\,mm behind the body centre in the stance pose; agreement to 0.2\,\%
confirms that the inertia aggregation, the MJCF generation and the contact model
are mutually consistent.

\subsection{Test environment}

\begin{decision}{One world description, two renderers}
\code{world\_spec.py} is consumed by the MJCF builder and by the RViz marker
publisher. What the operator sees in RViz is geometrically what the robot
collides with in MuJoCo; two hand-maintained files would drift, and the drift
would look like a perception bug.
\end{decision}

Two worlds ship. \code{flat} is an open-air proving ground, 24\,$\times$\,24\,m,
with eight independent lanes radiating from the spawn so that each locomotion
problem can be attempted on its own: a measured run with metre markers for
velocity tuning, a step ladder from 40 to 200\,mm, a stair complex (six
80\,mm steps up, a landing, six down, and a 130\,mm flight beside it), slopes
at 5, 10, 15 and 20$^\circ$, a 196-tile rough-terrain field, stepping stones,
a gap course widening from 100 to 280\,mm, and a 220\,mm balance beam with a
pole slalom. Gait work should not be debugged against furniture at the same
time, which is why it is separate from the house.

\code{house} is 12\,$\times$\,9\,m and contains every indoor locomotion
problem the robot must solve, each with enough approach room to be attempted in
isolation:

\begin{center}\small
\begin{tabular}{@{}lll@{}}
\toprule
\textbf{Feature} & \textbf{Dimension} & \textbf{Against} \\
\midrule
Doorways         & 0.90\,m, 20\,mm sill & body width 0.24\,m \\
Narrow gaps      & 0.45 / 0.43 / 0.60\,m & foot span 0.289\,m \\
Stairs           & 4 $\times$ 80\,mm; 3 $\times$ 120\,mm & realistic and stretch \\
Ramp             & 12$^\circ$ & friction limit at $\mu = 0.9$ \\
Under-table      & 0.36--0.46\,m clearance & stance height \RdStandHeight\,mm \\
Debris field     & mixed 40--110\,mm blocks & foothold selection \\
Movable props    & 4 free bodies & contact recovery \\
\bottomrule
\end{tabular}
\end{center}

\code{world\_check.py} measures every one of those clearances from the spec, and
\code{test\_world.py} fails if an edit turns a designed challenge into a wall or
an open field. It also rejects any overlapping geometry, because MuJoCo resolves
interpenetration by pushing bodies apart at the first step --- silently moving
the test environment out from under the experiment.

% =====================================================================
\section{Perception}

\archfig{perception_pipeline}{One configuration file drives both backends, so
the simulated camera has the real device's intrinsics and range.}{1.0}

The Yahboom NUWA HP60C is mounted on the body's front face on a bracket at
$z = 175$\,mm, pitched 0.26\,rad down.

\begin{decision}{The camera mount height is set by self-occlusion}
The front hip assemblies occupy roughly $x = 0.11 \ldots 0.26$\,m and reach
$z = +0.05$\,m, leaving only a 42\,mm slot between them on the centreline.
Rendering the simulated depth image from $z = 20$\,mm put \textbf{95\,\%} of the
frame on the robot's own legs at 74\,mm; $z = 90$\,mm still left 48\,\%. Clearing
the hip corner with the bottom edge of the frustum requires $z > 155$\,mm, so
the camera sits on a mast --- as it does on every production quadruped. At
175\,mm the self-view is 0\,\% and 98.5\,\% of the depth image is usable.
\end{decision}

\begin{decision}{The simulated camera loads its own MuJoCo model}
The physics instance lives in the control node's backend. Rather than reach into
it, the camera backend loads the same scene read-only and drives it from the
published robot state. The camera is then a separate node with a separate
lifecycle --- exactly as the real USB camera will be --- so the node graph does
not change on hardware, rendering never steals time from the
\RdControlRate\,Hz loop, and a camera crash cannot take the control loop down.
\end{decision}

Depth crosses the backend boundary as float32 metres with NaN for invalid,
because that representation can express ``no reading''; the conversion to the
ROS 16UC1-millimetre convention (where 0 means invalid) happens once, at the
publisher. Stereo range error grows as $z^2$, so the simulated camera applies
$\sigma_z = z^2 \sigma_d / (f b)$ --- about 4\,mm at 1\,m and 140\,mm at 6\,m ---
and the point cloud is filtered to \code{max\_usable\_range\_m} rather than the
datasheet maximum.

\begin{caveat}{Camera parameters are placeholders}
Resolution, field of view, range and baseline in
\code{config/nuwa\_hp60c.yaml} are plausible for a device of this class but have
\emph{not} been read from the datasheet; each is marked \code{[VERIFY]}. The
architecture does not depend on them --- they are configuration, and both
backends read the same file. How the HP60C delivers depth over USB (a second UVC
stream, a side-by-side frame, or only through the vendor SDK) is the one thing
that cannot be determined without the device, and is marked \code{HW-CHECK} in
\code{backends/nuwa\_hp60c.py}.
\end{caveat}

% =====================================================================
\section{Web GUI}

A single process holds an rclpy node and an aiohttp server; rclpy spins in a
background thread while aiohttp owns the asyncio loop. They meet at two
dictionaries: the ROS callbacks write the latest state, the web tasks read it.
No queue --- a GUI only ever wants the newest value, and buffering telemetry for
a browser that fell behind is how a monitoring tool ends up showing the past.

\begin{decision}{A purpose-built protocol rather than rosbridge}
rosbridge would re-serialise every field at the full publish rate and expose the
whole graph to the browser. This speaks a small versioned JSON protocol,
decimated to 20\,Hz, and exposes exactly the commands the GUI is allowed to
send. Less code, less bandwidth, and a much smaller surface from a page that can
drive a robot.
\end{decision}

The browser fetches \code{/api/config} and builds itself from it: adding,
removing or reordering a panel is a change to \code{config/web.yaml} alone.
Panels cover robot state, controller state, per-joint telemetry with tracking
error and load bars, foot contact and forces, safety and thermal state,
simulation state and real-time factor, an event log, camera streams, pose and
gait commands, and emergency stop.

Operator limits are configuration, not constants in JavaScript: maximum
commanded velocity, gait confirmation, and joint jogging --- which is disabled
by default because it bypasses the gait layer and is the easiest way to put a
leg where the body cannot support it. Video is a separate MJPEG endpoint so a
slow video consumer cannot stall telemetry, and the page falls back to showing
camera metadata if no JPEG encoder is installed.

% =====================================================================
\section{Migration to hardware}
\label{sec:migration}

\subsection{What changes}

\begin{center}\small
\begin{tabular}{@{}lll@{}}
\toprule
\textbf{Component} & \textbf{Simulation} & \textbf{Hardware} \\
\midrule
Joints  & \code{MujocoBackend} & \code{RobStride02Backend} \\
Camera  & \code{SimCameraBackend} & \code{NuwaHP60CBackend} \\
Base state & ground truth from the simulator & \code{StateEstimator} \\
Temperature & thermal model & reported, with the model as observer \\
Launch  & \code{backend:=mujoco} & \code{backend:=robstride02\_can} \\
\bottomrule
\end{tabular}
\end{center}

Nothing else: not a topic, a service, a frame, a message, a controller or a gain.

\subsection{Bring-up sequence}

\begin{enumerate}
\item \textbf{Verify the CAN protocol.} Power one motor alone; confirm
      \code{GET\_DEVICE\_ID}, then \code{ENABLE} and the arrival of feedback
      frames, then a 0.1\,N$\cdot$m command with $k_p = k_d = 0$ and check the
      direction of motion. A firmware revision with a different identifier
      layout changes exactly one file.
\item \textbf{Calibrate each joint.} \code{ros2 run robodog\_hardware
      calibrate\_joint} determines \code{direction} and \code{offset\_rad}
      against the convention in Section~\ref{sec:conventions} and prints a YAML
      fragment for review. Leave \code{enable\_on\_start: false} until signed off.
\item \textbf{Identify the estimates.} Rotor inertia from a free-spin
      deceleration test; joint friction and damping from a swing-down test;
      thermal constants from a stall at rated torque. Until then the margins in
      Section~\ref{sec:safety} carry the uncertainty.
\item \textbf{Bring up on a stand}, legs free: joint test sweeps, then named
      poses, then stance with the body supported.
\item \textbf{On the ground}, walk gait first --- it is statically stable and
      keeps three feet down.
\end{enumerate}

\subsection{Known gaps}
\label{sec:assumptions}

\begin{center}\small
\begin{tabular}{@{}lll@{}}
\toprule
\textbf{Item} & \textbf{Status} & \textbf{Resolution} \\
\midrule
CAN frame layout          & from the RS02 manual & verify against delivered firmware \\
Joint direction / offset  & unknown & \code{calibrate\_joint}, per joint \\
Rotor inertia             & estimate \RdRotorInertia & free-spin deceleration \\
Damping, Coulomb friction & estimate & swing-down test \\
Thermal $R_{\mathrm{th}}$, $C_{\mathrm{th}}$ & estimate & stall at rated torque \\
Temperature thresholds    & estimate & from the above, plus vendor guidance \\
Camera intrinsics, range  & placeholder & datasheet, then \code{camera\_calibration} \\
HP60C depth transport     & unknown & determine with the device \\
Joint limits              & design values & self-collision check on real geometry \\
Thigh/base clearance      & contact excluded & confirm in CAD; see \S\ref{sec:locomotion-status} \\
Battery model             & nominal voltage only & add a cell model and state of charge \\
\bottomrule
\end{tabular}
\end{center}

\subsection{Planned next steps}

The \code{ros2\_control} interface contract is already declared in
\code{urdf/ros2\_control.xacro} but not implemented: the runtime is Python.

\begin{decision}{Python control loop now, C++ \code{SystemInterface} later}
At \RdControlRate\,Hz the measured loop period is 2.49\,ms against a nominal
2.50\,ms, with worst-case jitter under 1\,ms --- adequate, and it buys a stack
that is fully testable and quick to change during architecture work. The
declared \code{ros2\_control} block is the formal statement of the interfaces a
C++ \code{SystemInterface} must provide, and is the migration path once jitter
or CPU headroom demands it. Note that the command interface set there is the
RS02's native impedance command, not a single position interface, for the reason
in Section~\ref{sec:actuator}.
\end{decision}

% =====================================================================
\section{Verification}

\begin{center}\small
\begin{tabular}{@{}lrl@{}}
\toprule
\textbf{Suite} & \textbf{Tests} & \textbf{What it protects} \\
\midrule
\code{robodog\_description} & 16 & URDF structure, conventions, limits, inertia validity \\
\code{robodog\_hardware}    & 52 & CAN codec, backend contract, thermal, bus budget \\
\code{robodog\_control}     & 73 & IK, Jacobian, force signs, safety, balance, gait continuity \\
\code{robodog\_sim}         & 59 & MJCF vs URDF consistency, world geometry, physics, locomotion \\
\code{robodog\_perception}  & 17 & intrinsics, deprojection, self-occlusion, noise model \\
\midrule
\textbf{Total}              & \textbf{217} & all pass \\
\bottomrule
\end{tabular}
\end{center}

The most valuable are the cross-checks: the MJCF's per-link masses, joint
origins, axes and ranges are asserted against the URDF, so the two descriptions
cannot drift. Seven tests exist specifically because a bug got past a weaker
one: the gravity-torque sign, the idle-command rate-limiter seed, the camera
self-occlusion fraction, the Euler-convention mismatch that left the rendered
legs beside the body while every physics test passed, the swing arc's
discontinuity at lift-off, the self-collision that welded the front hips, and
the gait's thermal load.

The last three are worth noting together. The thermal test spent most of the
project as an expected failure at 173\,\% of continuous, and it did not turn
green by being tuned --- it turned green when the two bugs above it were
fixed, and it now reads 90\,\%. An expected-failure test is a good way to hold
a known gap where a counter will find it, but it says nothing about
\emph{why} the gap exists: the reason sat two layers below where the test
pointed, in the simulation model.

\subsection{Analysis tools}

Everything quoted in this document is measurable, and the measurement is in the
repository rather than in a notebook someone once ran.

\begin{center}\small
\begin{tabular}{@{}ll@{}}
\toprule
\textbf{Tool} & \textbf{Answers} \\
\midrule
\code{cad\_to\_model.py}      & CAD $\to$ canonical model; prints the FK residual \\
\code{prepare\_meshes.py}     & mesh decimation, 41\,MB $\to$ 2.3\,MB \\
\code{torque\_report.py}      & peak and RMS torque per gait, and the winding temperature \\
\code{capability\_report.py}  & jump, payload, step-up, slope and power envelope \\
\code{design\_study.py}       & torque against leg length and stance height \\
\code{optimise\_stance.py}    & the stance-pose search behind the named poses \\
\code{make\_diagrams.py}      & the ten figures, as Draw.io and PDF \\
\code{make\_torque\_figure.py} & why torque follows leg bend, not leg length \\
\code{make\_doc\_facts.py}    & every number in this document \\
\code{render\_previews.py}    & the MuJoCo preview renders \\
\bottomrule
\end{tabular}
\end{center}

\subsection{Capability envelope}
\label{sec:capability}

Measured with \code{capability\_report.py}, using deliberately crude open-loop
controllers so the result describes the machine rather than any particular
control scheme. This is a separate question from whether the current gait
works, which Section~\ref{sec:locomotion-status} answers.

\begin{center}\small
\begin{tabular}{@{}lll@{}}
\toprule
\textbf{Capability} & \textbf{Measured} & \textbf{Limited by} \\
\midrule
Walking, 0.5\,m/s  & $\approx$49\,W net of 754\,W rated & nothing; 7\,\% of budget \\
Stairs, 80\,mm     & needs 100\,mm lift, has 130\,mm & leg travel \\
Maximum step       & 110\,mm standing, 190\,mm from a crouch & leg travel \\
Vertical jump      & 277\,mm of flight, 308\,ms airborne & peak torque \\
Landing            & falls over & \emph{controller}, not hardware \\
Payload            & $\approx$5\,kg before stance exceeds continuous & continuous torque \\
Slope              & 30$^\circ$ at 55\,\% of continuous & friction beyond 35$^\circ$ \\
Peak power         & 2.19\,kW, 219\,W/kg & --- \\
\bottomrule
\end{tabular}
\end{center}

The useful number is the first one. Steady walking costs about 7\,\% of the
rated mechanical power, because almost all the torque a quadruped spends is
holding itself up, which does no work. That is why the gait's
\RdTorqueCont\,N$\cdot$m-exceeding draw in Section~\ref{sec:thermal-load} is
diagnostic rather than a sizing problem: the machine is not working hard, the
controller is wasting what it asks for.

\archfig{domain_model}{Domain model: the value types that cross the hardware
boundary and the entities that own them.}{0.94}

\vspace{4mm}
\begin{center}\small\color{rdgrey}
Generated from \code{robot\_parameters.yaml} (CAD \texttt{\RdCadSha}) and
\code{robstride02.yaml}.\\
Diagrams: \code{docs/diagrams/*.drawio}, rendered by \code{tools/make\_diagrams.py}.
\end{center}

\end{document}
